% Generated draft: requires mathematical and source review.
\documentclass[11pt]{article}
\usepackage[margin=25mm]{geometry}
\usepackage{fontspec}
\setmainfont{Latin Modern Roman}
\usepackage{amsmath,amssymb,mathtools,unicode-math}
\setmathfont{Latin Modern Math}
\usepackage{enumitem,needspace,xurl,hyperref}
\setlength{\parindent}{0pt}
\setlength{\parskip}{5pt}
\setlength{\emergencystretch}{3em}
\newcommand{\N}{\mathbb{N}}
\newcommand{\Z}{\mathbb{Z}}
\newcommand{\Q}{\mathbb{Q}}
\newcommand{\R}{\mathbb{R}}
\newcommand{\C}{\mathbb{C}}
\newcommand{\F}{\mathbb{F}}
\newcommand{\A}{\mathbb{A}}
\newcommand{\PP}{\mathbb{P}}
\newcommand{\E}{\mathbb{E}}
\newcommand{\eps}{\varepsilon}
\newcommand{\dd}{\,\mathrm{d}}
\begin{document}
\section*{UnsolvedMath ID: 701736}
\textbf{Source problem code:} catalogue entry
\subsection{Definitions and mathematical statement}
Let \(G\) be a finite group, and let \(\operatorname{Irr}(G)\) denote its set of irreducible complex characters. For \(\chi\in\operatorname{Irr}(G)\), its degree is \(\chi(1)\), and its kernel is
\[
\ker\chi=\{g\in G:\chi(g)=\chi(1)\}.
\]
Define its \emph{codegree} and the set of character codegrees of \(G\) by
\[
\operatorname{cod}(\chi)=\frac{|G:\ker\chi|}{\chi(1)},
\qquad
\operatorname{Cod}(G)=\{\operatorname{cod}(\chi):\chi\in\operatorname{Irr}(G)\}.
\]
\textbf{Conjecture (G. Qian).} For every finite group \(G\) and every positive integer \(m\), if \(G\) contains an element of order \(m\), then there exists \(\chi\in\operatorname{Irr}(G)\) such that
\[
m\mid\operatorname{cod}(\chi).
\]
Equivalently, under this hypothesis there exist \(\chi\in\operatorname{Irr}(G)\) and a positive integer \(k\) such that
\[
|G:\ker\chi|=km\chi(1).
\]
\textbf{Supplied qualifications.} The record reports that the conjecture has been proved when \(m\) is a prime power (G. Qian, 2011), when \(m\) is square-free (I. M. Isaacs, 2011), when \(G\) is solvable (G. Qian, 2021), when \(G\) is a symmetric or alternating group (E. Giannelli, 2024), when \(G\) is almost simple (S. Y. Madanha, 2023), and when \(\mathbf F(G)=1\) (Z. Akhlaghi, E. Pacifici, and L. Sanus, 2024). Here \(\mathbf F(G)\) is the Fitting subgroup, the largest nilpotent normal subgroup of \(G\); an almost simple group satisfies \(S\leq G\leq\operatorname{Aut}(S)\) for some nonabelian finite simple group \(S\). These special-case results and the general conjecture's open status are supplied source assessments, not independently verified conclusions.
\subsection{Short English statement}
The conjecture asks whether every element order in a finite group divides the codegree of some irreducible complex character of that group. The supplied record labels the general problem open and lists several proved special cases; those literature assessments require verification.
\subsection{Sources}
\begin{itemize}
\item Local UnsolvedMath registry, permanent ID \(701736\), exact source code \texttt{catalogue entry}, title \emph{Kourovka Notebook Problem 20.78}. The supplied record attributes the problem to G. Qian, with proposed year 2022, and labels it \texttt{open}. Its recorded update date is 2026-10-16; this is registry metadata, not independent literature verification.
\item \emph{The Kourovka Notebook (Unsolved Problems in Group Theory)}, Problem 20.78. The supplied background identifies the 20th issue (2022) and the 21st issue edition (2026), \texttt{arXiv:1401.0300v48}. The supplied collection description names E. I. Khukhro and V. D. Mazurov as editors.
\item G. Qian, \emph{Arch. Math.}, volume 97 (2011), pages 99--103. Supplied reference for the prime-power case.
\item I. M. Isaacs, \emph{Arch. Math.}, volume 97 (2011), pages 499--501. Supplied reference for the square-free case.
\item G. Qian, \emph{Bull. London Math. Soc.}, volume 53 (2021), pages 820--824. Supplied reference for solvable groups.
\item E. Giannelli, \emph{J. Algebra Appl.}, volume 23, number 9 (2024), article ID 2450144. Supplied reference for symmetric and alternating groups.
\item S. Y. Madanha, \emph{Commun. Algebra}, volume 51, number 7 (2023), pages 3143--3151. Supplied reference for almost simple groups.
\item Z. Akhlaghi, E. Pacifici, and L. Sanus, \emph{J. Algebra}, volume 644 (2024), pages 428--441. Supplied reference for groups with trivial Fitting subgroup.
\end{itemize}
\end{document}
